\documentclass{UVKABoA4}

\usepackage{graphicx} % Standardpaket zur Grafikeinbindung
\usepackage{amsmath,amssymb} % Erweiterung des Mathematik-Modus
\usepackage[colorinlistoftodos, german]{todonotes} % Option 'disable' entfernt alle ToDos
\usepackage[absolute,overlay]{textpos}
\usepackage{vmargin}          % Adjust margins in a simple way
\usepackage{tikz}

% Toggle the following two lines to switch between english and german layout
% \usepackage[english]{babel} % Neue deutsche Rechtschreibung und Silbentrennung.
\usepackage[ngerman]{babel} % Neue deutsche Rechtschreibung und Silbentrennung.

\usepackage[raiselinks=true,
            bookmarks=true,
            bookmarksopenlevel=1,
            bookmarksopen=true,
            bookmarksnumbered=true,
            hyperindex=true,
            plainpages=false,
            pdfpagelabels=true,
            pdfborder={0 0 0.5}]{hyperref}

\renewcommand{\labelitemi}{$\triangleright$} 

\makeatletter
\def\tagform@#1{\maketag@@@{\ignorespaces [#1]\unskip\@@italiccorr}}	% Anpassung der Formelnummerierung
\makeatother

\makeindex

\newcommand{\editora}{Abel~Gawel}
\newcommand{\editorb}{Anne~Günter}
\newcommand{\editorc}{Hamza~Hajj~Hassan}
\newcommand{\editord}{M\'{a}rton~Kelemen}
\newcommand{\editore}{Oier~Mees}
\newcommand{\editorf}{Tobias~Roth}
\newcommand{\editorg}{Dirk~Toewe}
\newcommand{\editorh}{Andreas~Weber}
\newcommand{\editori}{Julian~Wormer}

\newcommand{\editorx}{Editor~X}

\newcommand{\headtitle}{Projektpraktikum mit mobilen Robotern}
\newcommand{\floatingtitle}{Projektpraktikum mit mobilen Robotern}
\newcommand{\thesistype}{Ausarbeitung}

\newcommand{\advisorone}	{Dipl.--Inform.~Andreas~Hermann}
\newcommand{\advisortwo}	{Dipl.--Inform.~Jan~Oberländer~(M.\,Sc.)}
\newcommand{\advisorthree}	{Dipl.--Inform.~Zhixhing~Xue}

\newcommand{\timestart}{Oktober 2011}
\newcommand{\releaseyear}{2012}
\newcommand{\timeend}{April \releaseyear}
\newcommand{\releasemonth}{im April \releaseyear}


% Names of the organizations
\newcommand{\department}{\iflanguage{english}{Department of Computer Science}
                                             {Fakultät für Informatik}}
\newcommand{\institute}{\iflanguage{english}{Institute for Anthropomatics}
                                            {Institut für Anthropomatik}}
\newcommand{\fzidepartment}{Abteilung Technisch Kognitive Assistenzsysteme}
\newcommand{\fziname}{\iflanguage{english}{FZI Research Center for Information Technology}
                                          {FZI Forschungszentrum Informatik}}

% A graphic best describing your thesis
\newcommand{\titlefig}{\includegraphics[width=0.4\textwidth]{graphics/FZI-Logo}}

\begin{document}

\listoftodos

\pagenumbering{roman}
%% titlepage.tex
%%

% coordinates for the bg shape on the titlepage
\newcommand{\diameter}{20}
\newcommand{\xone}{-25}
\newcommand{\xtwo}{150}
\newcommand{\yone}{25}
\newcommand{\ytwo}{-243}

\begin{titlepage}
% bg shape
\begin{tikzpicture}[overlay]
\draw[color=gray]  
 		 (\xone mm, \yone mm)
  -- (\xtwo mm, \yone mm)
 arc (90:0:\diameter pt) 
  -- (\xtwo mm + \diameter pt , \ytwo mm) 
	-- (\xone mm + \diameter pt , \ytwo mm)
 arc (270:180:\diameter pt)
	-- (\xone mm, \yone mm);
\end{tikzpicture}
  \begin{textblock}{10}[0,0](3.35,2.55)
	  \iflanguage{english}  {\includegraphics[width=.3\textwidth]{graphics/kitlogo_en_rgb}}
                          {\includegraphics[width=.3\textwidth]{graphics/kitlogo_de_rgb}}
  \end{textblock}
	\changefont{phv}{m}{n}	% helvetica	
	\vspace*{0.5cm}
	\begin{center}
		\Huge{\headtitle}
		\vspace*{1cm}\\
		\Large{
			\iflanguage{english}{Diploma Thesis of}			
												  {Ausarbeitung\\von}
		}\\
		\vspace*{1cm}
		\Large{\editora}	\\
		\Large{\editorb}	\\
		\Large{\editorc}	\\
		\Large{\editord}	\\
		\Large{\editore}	\\
		\Large{\editorf}		\\
		\Large{\editorg}	\\
		\Large{\editorh}	\\
		\Large{\editori}		\\
		\vspace*{1.5cm}
		\Large{
			\department\\ \institute\\ \iflanguage{english}{and}{und}\\ \fziname
		}
	\end{center}
	\vspace*{0.5cm}
\Large{
\begin{center}
\begin{tabular}[ht]{l c l}
 % \iflanguage{english}{Reviewer}{Erstgutachter}: & \hfill  & \reviewerone\\
 % \iflanguage{english}{Second reviewer}{Zweitgutachter}: & \hfill  & \reviewertwo\\
  \iflanguage{english}{Advisor}{Betreuende Mitarbeiter}: & \hfill	& \advisorone		\\
  																						&			& \advisortwo		\\
  																						&			& \advisorthree	\\
  % \iflanguage{english}{Second advisor}{Zweiter betreuender Mitarbeiter}: & \hfill  & \advisortwo\\
\end{tabular}
\end{center}
}


\vspace{1cm}
\begin{center}
\large{\iflanguage{english}{Research Period}{Bearbeitungszeit}: \timestart \hspace*{0.25cm} -- \hspace*{0.25cm} \timeend}
\end{center}


\begin{textblock}{10}[0,0](4,16.8)
\tiny{ 
	\iflanguage{english}
		{KIT -- University of the State of Baden-Wuerttemberg and National Research Center of the Helmholtz Association}
		{KIT -- Universität des Landes Baden-Württemberg und nationales Forschungszentrum in der Helmholtz-Gemeinschaft}
}
\end{textblock}

\begin{textblock}{10}[0,0](14,16.75)
\large{
	\textbf{www.kit.edu} 
}
\end{textblock}

\end{titlepage}

\pagestyle{empty}
\begingroup
\changefont{phv}{m}{n}
\renewcommand{\baselinestretch}{1}

\begin{titlepage}
  \parindent 0em 
  \LARGE\textbf{\floatingtitle}
  \vspace{24pt}\\
  \iflanguage{english}{by}{von}\\
  \myname \\[1.5cm]
  \begin{center}
  \titlefig\\[1cm]
  \end{center}
  \vfill
  \begin{minipage}{0.4\textwidth}
    \begin{flushleft} \large
      \LARGE\textbf{\thesistype}\\
      \LARGE\releasemonth
    \end{flushleft}
  \end{minipage}
  \begin{minipage}{0.6\textwidth}
    \begin{flushright}
      \includegraphics[width=1.5cm]{graphics/FZI-Logo}
    \end{flushright}
  \end{minipage}
  \newpage
  \small \thesistype, FZI\\
  \department, \releaseyear\\
  Gutachter: \reviewerone, \reviewertwo
  \vfill
  \fzidepartment\\
  \fziname
  \newpage
\end{titlepage}

\endgroup

% \thispagestyle{empty}

\begin{center}
  \vspace*{\stretch{1}}
  \bfseries
  \LARGE
  \headtitle\\[0.45cm]
  \vfill
  \mdseries
  \huge
  \thesistype\\[0.8cm]
  \large
  Abteilung Technisch Kognitive Assistenzsysteme\\[0.1cm]
  Forschungszentrum Informatik\\[0.1cm]
  an der\\[0.1cm]
  Universität Karlsruhe (TH)\\
  \vfill
  von\\[0.3cm]
  \LARGE
  \myname\\
  \vfill
  \normalsize
  \begin{tabular}{p{3.5cm} p{0.4cm} p{5.5cm}}
    \textbf{Tag der Ausgabe} & : &  \timestart\\
    \textbf{Tag der Abgabe}  & : &  \closingdate\\
    & & \\
    & & \\
    & & \\
    \textbf{Betreuer} & : & \advisor \\
    \textbf{Referent} & : & Prof.~Dr.--Ing.~J.~M.~Zöllner \\
    \textbf{Koreferent} & : & Prof.~Dr.--Ing.~R.~Dillmann
  \end{tabular}
\end{center}

\parindent 1em
% \affirmation

Ich versichere wahrheitsgemäß, die Arbeit selbstständig angefertigt, alle benutzten Hilfsmittel vollständig und genau angegeben und alles kenntlich gemacht zu haben, was aus Arbeiten anderer unverändert oder mit Abänderungen entnommen wurde.

\vspace{2cm}
\begin{flushright}\noindent
    Karlsruhe,\hfill {\it \myname}\\
    \releasemonth \hfill { }
\end{flushright}
% \preface

Diese \LaTeX{}-Vorlage, basierend auf \emph{KOMA}-Script-Klasse\index{KOMA-Script} scrbook\index{scrbook}, wurde erstellt, um Veröffentlichungen im Universitätsverlag\index{Universitätsverlag} Karlsruhe in einem einheitlichen Layout zu ermöglichen. Die Dokumentklasse\index{Dokumentklasse} \emph{UVKABook}\index{UVKABook} soll dabei unverändert bleiben. Sollten zusätzliche Pakete geladen werden, muss sichergestellt sein, dass keine Veränderungen am bestehenden Layout erfolgen. Der Inhalt des Manuskripts wird im Unterordner \emph{content} abgelegt und in das Hauptdokument \emph{book.tex} eingebunden. Bestehende Abschnitte etwa für das Vorwort\index{Vorwort} oder die Danksagung\index{Danksagung} können einkommentiert werden. Unter Umständen kann es notwendig sein, dass einzelne Pakete nachinstalliert werden (beispielweise das \emph{KOMA-Script} oder \emph{natbib}\index{natbib}). Das Literaturverzeichnis\index{Literaturverzeichnis} verwendet \emph{bibtex}. Bei Rückfragen wenden Sie sich bitte an den Universitätsverlag Karlsruhe.

\vspace{1cm}
\begin{flushright}\noindent
    Karlsruhe,\hfill {\it \myname}\\
    \releasemonth \hfill { }
\end{flushright}
% 
\abstract

Agnisim ad te diamcon sectet, voloborper aciduis nonsequamet, sequisciduis nonum irilit ea feugait, sequat vero er ip exero odipit lobore magnit nos nibh ex esed exerciduisi tat. Consed et  inci tet aci et atem in ut velenim vent ullum exeril ut laore molorpe riliscidunt la con ulputat ad digna feum quismolesto dolortisl ea feuguercil delit praessim ver ad tem ipit, velit, conse
enis nonsed mod tincipit augue veniametuer aut vulput nonsed molesse miniatuercin hendit in henibh etumsan velessi.
It verilluptat adiatie vel ullum il dolorem zzriuscing et pratinit wis acilismodit adions ad magna autat ulput adip eum do eu faci tem ver ipis at augue ming erit wiscillan hent iliquisl ullam ecte delent ent adiam zzrit aliqui tie volorem inibh essequip enis nullaore vent ametum vullam, velese vullametue verit, sit ex er il ute dolortio commy nos el utat. 
Nons dolore voloborper sed magnit enisisi. Onsenit digna facip ex eraestin ulla consequatie do cons nulla feugiat. Bor irilla feuguer cipsum er autpat. Duis ad minciliquat. Er sis autat ut wis dolorper adion hendre ver sequate molore etuerci exeros nonulla consectet lorpera essequiscip ectet, suscil el utpat luptatum velis nulla...
% 
\ack

Die vorliegende Arbeit entstand während meiner Tätigkeit als wissenschaftlicher Mitarbeiter am Institut für Formalerschließung und Regelwerkskunde der Universität Karlsruhe (TH). Herrn Prof. Dr.-Ing. habil. A. Mustermann danke ich besonders für die Anregung zu dieser Arbeit, die wissenschaftliche Förderung, die stets vorhandene Diskussionsbereitschaft und für die Übernahme des Hauptreferates.

Für die freundliche Übernahme des Korreferates gebührt mein ganz besonderer Dank Herrn Prof. Dr.-Ing. B. Musterfrau vom Institut für Formalerschließung und Regelwerkskunde. Dem Vorsitzenden des Prüfungsausschusses, Herrn Prof. Dr.-Ing. C. Mustermann gilt ebenfalls mein Dank für die kritische Durchsicht des Manuskriptes. Meinen Eltern, die mir durch ihre Redlichkeit und Strebsamkeit immer Vorbild waren, die mir Heimat und Geborgenheit aber auch Freiheit und Interesse für Neues gaben, danke ich von tiefstem Herzen, ebenso gilt der Dank meinen Geschwistern. Für die Geduld und Entbehrungen, die meine Frau Berta und mein Sohn Anton während der letzten Monate aufbrachten und für ihre unendliche Motivation und ihr Verständnis, das sie mir entgegenbrachten, danke ich zutiefst.

\vspace{1cm}
\begin{flushright}\noindent
    Karlsruhe,\hfill {\it \myname}\\
    \releasemonth \hfill { }
\end{flushright}

\begingroup
\changefont{phv}{m}{n}
\tableofcontents
\endgroup

\mainmatter
\renewcommand{\chapterpagestyle}{plain}
\pagestyle{scrheadings}
\pagenumbering{arabic}


\chapter{Einleitung}

\todo[inline]{Struktur erstmal übernommen von WS09/10\ldots}

\section{Aufgabenstellung}
\todo[inline]{Schreiben}
Das hat der Tobias sehr schön gemacht, finde ich.
Lorem ipsum dolor sit amet, consetetur sadipscing elitr, sed diam nonumy eirmod tempor invidunt ut labore et dolore magna aliquyam erat, sed diam voluptua.
At vero eos et accusam et justo duo dolores et ea rebum.
Stet clita kasd gubergren, no sea takimata sanctus est Lorem ipsum dolor sit amet.
Lorem ipsum dolor sit amet, consetetur sadipscing elitr, sed diam nonumy eirmod tempor invidunt ut labore et dolore magna aliquyam erat, sed diam voluptua.
At vero eos et accusam et justo duo dolores et ea rebum.
Stet clita kasd gubergren, no sea takimata sanctus est Lorem ipsum dolor sit amet.


\section{Motivation}
\todo[inline]{Schreiben}
Lorem ipsum dolor sit amet, consetetur sadipscing elitr, sed diam nonumy eirmod tempor invidunt ut labore et dolore magna aliquyam erat, sed diam voluptua.
At vero eos et accusam et justo duo dolores et ea rebum.
Stet clita kasd gubergren, no sea takimata sanctus est Lorem ipsum dolor sit amet.
Lorem ipsum dolor sit amet, consetetur sadipscing elitr, sed diam nonumy eirmod tempor invidunt ut labore et dolore magna aliquyam erat, sed diam voluptua.
At vero eos et accusam et justo duo dolores et ea rebum.
Stet clita kasd gubergren, no sea takimata sanctus est Lorem ipsum dolor sit amet.


\section{Stand der Technik}
\todo[inline]{Schreiben}
Lorem ipsum dolor sit amet, consetetur sadipscing elitr, sed diam nonumy eirmod tempor invidunt ut labore et dolore magna aliquyam erat, sed diam voluptua.
At vero eos et accusam et justo duo dolores et ea rebum.
Stet clita kasd gubergren, no sea takimata sanctus est Lorem ipsum dolor sit amet.
Lorem ipsum dolor sit amet, consetetur sadipscing elitr, sed diam nonumy eirmod tempor invidunt ut labore et dolore magna aliquyam erat, sed diam voluptua.
At vero eos et accusam et justo duo dolores et ea rebum.
Stet clita kasd gubergren, no sea takimata sanctus est Lorem ipsum dolor sit amet.


\section{MCA2}
\todo[inline]{Anne? Schreiben}
Lorem ipsum dolor sit amet, consetetur sadipscing elitr, sed diam nonumy eirmod tempor invidunt ut labore et dolore magna aliquyam erat, sed diam voluptua.
At vero eos et accusam et justo duo dolores et ea rebum.
Stet clita kasd gubergren, no sea takimata sanctus est Lorem ipsum dolor sit amet.
Lorem ipsum dolor sit amet, consetetur sadipscing elitr, sed diam nonumy eirmod tempor invidunt ut labore et dolore magna aliquyam erat, sed diam voluptua.
At vero eos et accusam et justo duo dolores et ea rebum.
Stet clita kasd gubergren, no sea takimata sanctus est Lorem ipsum dolor sit amet.


\section{Systemübersicht}

\subsection{Hardware}
\todo[inline]{Schreiben}
Lorem ipsum dolor sit amet, consetetur sadipscing elitr, sed diam nonumy eirmod tempor invidunt ut labore et dolore magna aliquyam erat, sed diam voluptua.
At vero eos et accusam et justo duo dolores et ea rebum.
Stet clita kasd gubergren, no sea takimata sanctus est Lorem ipsum dolor sit amet.
Lorem ipsum dolor sit amet, consetetur sadipscing elitr, sed diam nonumy eirmod tempor invidunt ut labore et dolore magna aliquyam erat, sed diam voluptua.
At vero eos et accusam et justo duo dolores et ea rebum.
Stet clita kasd gubergren, no sea takimata sanctus est Lorem ipsum dolor sit amet.

\subsection{Software}
\todo[inline]{Schreiben}
Lorem ipsum dolor sit amet, consetetur sadipscing elitr, sed diam nonumy eirmod tempor invidunt ut labore et dolore magna aliquyam erat, sed diam voluptua.
At vero eos et accusam et justo duo dolores et ea rebum.
Stet clita kasd gubergren, no sea takimata sanctus est Lorem ipsum dolor sit amet.
Lorem ipsum dolor sit amet, consetetur sadipscing elitr, sed diam nonumy eirmod tempor invidunt ut labore et dolore magna aliquyam erat, sed diam voluptua.
At vero eos et accusam et justo duo dolores et ea rebum.
Stet clita kasd gubergren, no sea takimata sanctus est Lorem ipsum dolor sit amet.



\chapter{Grundlagen}

\section{MCA2}
\todo[inline]{Anne: Schreiben}


\section{Odete}
\todo[inline]{Dirk: Schreiben}


\section{SegwayOmni}
\todo[inline]{Dirk: Schreiben}


% #############################################################################
%															Bahnplanung und Steuerung
% #############################################################################
\chapter{Bahnplanung und Steuerung}
\label{bahnplanung_steuerung_cha}

% ********************************************************************************
% 										Aufgabenstellung
% ********************************************************************************
\section{Aufgabenstellung}
\label{bahnplanung_aufgabenstellung_sec}
\authorsection{\editoroier}


Was die Aufgabestellung für die Bahnplanung und Steuerung betrifft, so kann man das in Abstraktionsschichten einteilen:
LowLevel-Ansteuerung, Inverse Kinematik, Bahnplanung und Kollisionsvermeidung.
Außerdem sollte man das alles auch Simulieren können.

Für die LowLevel Ansteuerung ist die plattformspezifische Ansteuerung und Simulation notwendig. Außerdem sollte man auch die Pose aus Odometriedaten bestimmen können.

Die inverse Kinematik dient dazu um aus einer gegebene Pose, eine oder mehrere, bis hin zu unendlich vielen Roboterkonfigurationen zu berechnen. Letzteres passiert bei wenn es Raumpunkte gibt die zu Singularitäten führen. Das heißt,  dass zum Beispiel bei einem 6 Achsigen Roboter, zwei Achsen kollinear werden. Im Falle einer mobilen Plattform, würde man als inverse Kinematik die Berechnung einer Radstellung zu einer Pose (x, y, yaw) betrachten. Die Bestimmung der Geschwindigkeit wird danach durch inverse Odometrie gemacht. Zu lösende Aufgaben in Bereich der inverse Kinematik waren: die Umrechnung der Trajektorie in Geschwindigkeiten, die Limitierung der Beschleunigung des Roboters und das Limitieren der Geschwindigkeiten unter Beachtung des Protective Field \ref{sec:vpf}.

Was die Bahnplanung betrifft, gibt es mehrere Aufgaben. Als erstes müssen die Menschendaten, sprich, die Daten von der Kinect Kamera verarbeitet werden. Danach soll die Trajektorie auf der der Roboter fahren soll bestimmt werden, welches nebenbei kollisionsfrei sein sollte. Als letztes sollte das speichern der abgefahrenen Trajektorie und das Laden der abgefahrene Trajektorie implementiert werden.


%Als Ausgangssituation hatte man die Grundlagen der Odete für die LowLevel-Ansteurung und die inverse Kinematik. Sonst waren, Bahnplanung/Steuerung und Kollisionsvermeidung zunächst nicht vorhanden. Eine Simulationsumgebung war vorhanden, es fehlte aber noch die Programmierung der Gui .
%
%Nach dem Bottom-Up Prinzip hat man zunächst, die LowLevel-Ansteuerung von der Odete angepasst und mittels Simulation getestet.

%Das \textit{DriveAlongPath} Modul, welches sich mit der inverse Kinematik beschäftigt, dient um aus eine Trajektorie aus 3 Punkte die nötige Koordinaten für \textit{OmniDrive} zu berechnen.
%Falls es sich bspw. um ein Industrieroboter handeln würde, dann würde die inverse Kinematik die nötige Gelenkwinkeleinstellungen berechnen, um mit dem Greifer eine neue Position zu erreichen.
%Danach würde die inverse Dynamik die entsprechende Motor-Kräfte und Momente berechnen.
%
%Die Bahnplanung berechnet Trajektoriepunkte auf Basis der Mensch-Position und Geschwindigkeit, welche von der Kinect-Kamera geliefert werden.
%Außerdem werden hier diverse Transformationen von Welt zu Roboter bzw. Bild-Koordinaten und zurück gemacht.
%Nebenbei, wird der Pfad gespeichert damit der Roboter es später alleine abfahren kann.
%
%Zum Schluß wird eine ständige Überprüfung der geplante Trajektorie gemacht, um Kollisionen zu vermeiden.
%Um Hindernisse zu umfahren, werden der A*-Algorithmus und ein Potentialfeld eingesetzt.
%Letzteres wirkt, dass es \glqq teuer\grqq \space ist in der nähe von Hindernisse zu fahren.



% ********************************************************************************
% 										Grundlagen Bahnplanung
% ********************************************************************************
\section{Grundlagen der Bahnplanung und Steuerung}


\subsection{Steuerung}
\label{bahnplanung_steuerung_sec}

Weil die gefundene Trajektorie normalerweise nicht glatt ist werden für die Steuerung unterschiedliche Interpolationsarten verwendet. Die Bahnsteuerung kann in Kartesischen Raum oder im Konfigurationsraum stattfinden. In Kartesischen Raum ist es näher an der zu lösenden Aufgabe, man muss dann aber für jeden Punkt die inverse Kinematik berechnen. Die Steuerung in Konfigurationsraum hingegen liegt näher an den Gelenken, dafür ist die Formulierung der Trajektorie umständlicher. Vor und Nachteile beider Optionen kann man auf der Tabelle \ref{fig:steurungProContra} sehen.


\begin{table}[h]
		\centering
		\begin{tabular}{| p{7cm} | p{7cm}|}
\hline
\textbf{Kartesischer Raum} & \textbf{Konfigurationsraum}\\
\hline
+ Bahn einfacher zu formulieren & + Ansteuerung der Gelenke ist
einfacher\\
+ Interpolation ist einfacher & + Trajektorie ist eindeutig und
berücksichtigt Grenzen\\

- Inverse Kinematik ist für jeden Trajektorienpunkt zu lösen & - Interpolation für mehrere
Gelenke\\
- Geplante Trajektorie nicht immer ausführbar & - Formulierung der Trajektorie umständlicher\\
\hline
		\end{tabular}
		\caption{\label{fig:steurungProContra} Vor und Nachteile der Steuerung im Kartesischen und Konfigurationsraum Quelle: \citep{rob1}}
		\end{table}	
Es gibt unterschiedliche Interpolationsarten:  Zirkular, Linear und Spline-Bahn oder auch approximative Bahnsteuerung.
Bei der Linearinterpolation wird einfach zwischen zwei Teiltrajektorien interpoliert wie man auf der Abbildung \ref{fig:linearinterpolation} sehen kann.
\begin{figure}[h]
\center
\includegraphics[scale=0.35]{graphics/linearinterpolation.png}
\caption{\label{fig:linearinterpolation} Beispiel der Linearinterpolation Quelle: \citep{rob1}}
\end{figure}
 Die Zirkularinterpolation schafft hingegen kreisförmige Verfahrwege zwischen zwei Punkte. Wenn man hingegen  Segmentweise interpoliert, werden die Endbedingungen der Teiltrajektorie i-1 und die Anfangsbedingungen der Teiltrajektorie i aneinander angepasst, so dass die Teiltrajektorien durch Splines beschrieben werden wie man auf der Abbildung \ref{fig:segmentinterpolation} sehen kann.
 \begin{figure}[h]
\center
\includegraphics[scale=0.35]{graphics/segmentinterpolation.png}
\caption{\label{fig:segmentinterpolation} Beispiel der Segmentweise Bahninterpolation Quelle: \citep{rob1}}
\end{figure}
 
 
 Bei der approximative Bahnsteuerung werden anders als bei der Bahninterpolation, nicht alle Kontrollpunkte der Trajektorie befahren, sondern die Kontrollpunkte beeinflussen den Bahnverlauf und werden approximiert. Neben Bernsteinpolynome werden hier besonders Überschleifen benutzt um sanft Kurven zu fahren wie man auf der Abbildung \ref{fig:ueberschleifeninterpolation} sieht. Letzteres wurde auch in Praktikum eingesetzt.
 \begin{figure}[h]
\center
\includegraphics[scale=0.35]{graphics/ueberschleifeninterpolation.png}
\caption{\label{fig:ueberschleifeninterpolation} Beispiel der approximierte Bahnsteuerung mit Überschleifen Quelle: \citep{rob1}}
\end{figure}




\label{bahnplanung_grundlagen_sec}
\authorsection{\editoroier}
\subsection{Bahnplanung}

Bahnplanung befaßt sich mit dem Problem, eine kollisionsfreie Bahn zu finden die  ein Robotersystem
von der aktuellen Position in die Zielposition überführt. Damit gehört es zu einer der Basisprobleme der Robotik. Die Verfahren kann man nach unterschiedliche Gesichtspunkte klassifizieren \cite{rob1}.

Nach Robotertyp:
\begin{itemize}
\item Bahnplanung für Manipulatoren (Schweißen bei Industrierobotern z.B.)
\item Bahnplanung für mobile Roboter
\item Bahnplanung für Laufmaschinen und anthropomorphe Systeme
\item Greif- und Montageplanung
\end{itemize}

Nach dem Zustandsraum wo die Planung stattfindet:
\begin{itemize}
\item Gelenkwinkelzustandsraum (Konfigurationsraum, der Raum alle möglichen Gelenkwinkelkonfigurationen des Roboters )
\item Euklidischer Raum (Arbeitsraum, 3D / 6D)
\item Sensorzustandsraum
\item Objektzustandsraum
\end{itemize}

Nach Vollständigkeit:
\begin{itemize}
\item Vollständige Verfahren (liefern immer eine Korrekte Lösung und können ermitteln ob keine Lösung existiert)
\item Probabilistisch Vollständige Verfahren (Falls eine Lösung existiert konvergiert die Wahrscheinlichkeit, dass eine Lösung gefunden wird, bei fortschreitender Zeit gegen 1. Existiert keine Lösung, kann dies nicht ermittelt werden
)
\end{itemize}

Nach Navigationsart:
\begin{itemize}
\item Globale Pfadplanung bedeutet, eine Trajektorie zum Ziel zu planen wenn dem Roboter eine Karte der Umgebung bevorliegt und er sich dort lokalisieren kann.
\item Lokale Pfadplanung wird eingesetzt um Hindernisse zu umfahren und wird bei mobilen Robotern meist in Konbination mit einem globalen Planer benutzt. 
\end{itemize}

Weil die explizite Berechnung des kollisionsfreien Konfigurationsraum bei Roboter mit mehr Freiheitsgrade als 3 sehr teuer ist, benutzt man oft den impliziten Konfigurationsraum. D. h., dass die Bahnplanung in Konfigurationsraum stattfindet, aber die Kollisionen und Abstände in Arbeitsraum berechnet werden und ins diskretisierte Konfigurationsraum transformiert \citep{innoKonz}. So wird der Konfigurationsraum während der Suche exploriert.

Der Ablauf eines Planungsverfahren ist normalerweise in zwei geteilt\cite{Russell2003}.
\begin{enumerate}
\item Eine Roadmap oder Graph aufbauen mittels \gls{Voronoi}, \gls{Sichtgraphen}, \gls{Zellenzerlegung} oder zufällige Knoten je nach Sampling Strategie(für \gls{PRM} Verfahren\cite{Thrun2005} z.B.).
\item Im Graph ein Weg von Startknoten zum Zielknoten finden mittels A*-Algorithmus oder andere Suchalgorithmen für Graphen.
\end{enumerate}

\subsubsection{Lokale Planer}
\label{bahnplanung_lokale_planer_sec}

\todo[inline]{Drei gängige Planer vorstellen, state of the art am Ende.}




% ********************************************************************************
% 										Umsetzung
% ********************************************************************************
\section{Umsetzung}
\label{bahnplanung_umsetzung_sec}


% -----------------------------------------------------------------------------
%													LowLevel-Ansteuerung
\subsection{Motoransteuerung}
\label{bahnplanung_motoansteuerung_subsec}
\authorsection{\editoroier}

\todo[inline]{Oier: Schreiben}
\todo[inline]{Was es schon gab im Motoransteuerungsteil (aus Odete/segway).
SegwayOmniHal Modul gab es schon, wir haben das Modul ergänzt um die plattformspezifische Simulation Segway500RMPSimulation. Eingang sind Geschwindigkeiten und Ausgang sind Plattformspezifische Geschwindigkeiten)}.

% evtl. auch so ein Diagramm wie Abb. 4.4 von WS10/11?
%		bzw. Abb. 2.3 von WS09/10


% -----------------------------------------------------------------------------
%													Inverse Kinematik
\subsection{Inverse Kinematik}
\label{inverse_kinematik_subsec}
\authorsection{\editorjulian}

\todo[inline]{Julian: Schreiben}


% -----------------------------------------------------------------------------
%													Bahnplanung
\subsection{Bahnplanung}
\label{bahnplanung_subsec}
\authorsection{\editortobias}


% \begin{itemize}
% 	\item Kurzer Grobüberblick über die Behaviours-Gruppe, hier werden Bahnplanungs-Aufgaben bearbeitet\todoprivate{ref zu behaviours-Bild}
% 	\begin{itemize}
% 		\item Berechnung von Trajektorien-Punkten auf Basis der Mensch-Position und Geschwindigkeit (Kinect)
% 		\item Schätzung der zukünftigen Position des Menschen
% 		\item Transformation von Welt- in Roboter- bzw. Bildkoordinaten und zurück\todoprivate{evtl. auch weg lassen}
% 		\item Abspeichern von Pfaden
% 		\item Abfahren von Pfaden
% 	\end{itemize}
% 	\item Module werden im Folgenden kurz beschrieben
% 	\item Kollisionsvermeidung als eigener Abschnitt
% \end{itemize}

Die Aufgaben der Bahnplanung werden in der \lstinline{MCA2}-Gruppe \lstinline{SegwayOmniBehaviours} bearbeitet.\todoprivate{ref zu behaviours-Bild}
Im Einzelnen sind dies:
\begin{itemize}
	\item Berechnung von Trajektorien-Punkten auf Basis der Mensch-Position und Geschwindigkeit: \lstinline{GenerateTrajectory}
	\item Schätzung der zukünftigen Position des Menschen: \lstinline{HumanPosition}
	\item Abspeichern von Pfaden: \lstinline{StoreRoboterPositionsToFile}
	\item Abfahren von Pfaden: \lstinline{LoadRoboterPositionsFromFile}
	\item Kollisionsvermeidung: \lstinline{ObstacleAvoidance}
\end{itemize}

Die Module innerhalb der Gruppe werden im Folgenden beschrieben.
Die Kollisionsvermeidung wird in einem gesonderten Abschnitt (s. Abs. \ref{kollisionsvermeidung_subsec}) behandelt.

\missingfigure[figwidth=0.7\textwidth]{Behaviours}


\subsubsection{Position des Menschen}

% \begin{itemize}
%     \item bekommt aktuelle Position und Geschwindigkeit des Menschen von Kinect
%     \item schätzt die zukünftige Position des Menschen anhand dieser Daten
%     \item gibt aktuelle und geschätzte Position des Menschen aus
%     \item Zeit-Parameter für die Schätzung: $t_{est}$
%     \item Berechnung nach der Formel: $s_{est} = s_{act} + v \cdot t_{est}$
% \end{itemize}

Das Modul \lstinline{HumaPosition} bekommt als Parameter die aktuelle Position sowie die aktuelle Geschwindigkeit der vom Roboter verfolgten Person übergeben.
Anhand dieser Daten wird die zukünftige Position des Menschen geschätzt.
Als Resultat gibt das Modul die geschätzte zukünftige Position des Menschen aus.
Die aktuelle Position wird ebenfalls mit ausgegeben.

Für die Berechnung der zukünftigen Position wird die Formel
\begin{equation}
	s_{est} = s_{act} + v \cdot t_{est}
\end{equation}
angewendet.
$s_{est}$ ist dabei die geschätzte zukünftige Position, $s_{act}$ beschreibt die aktuelle Position und $v$ die Bewegungsgeschwindigkeit des Menschen.
$t_{est}$ ist ein Parameter für die Zeitspanne, die in die Zukunft prädiziert werden soll.
Dieser Parameter wird aus einer externen Konfigurationsdatei geladen.


\subsubsection{Berechnung der Trajektorie}

Die Berechnung der Trajektorie erfolgt im Modul \lstinline{GenerateTrajectory}.
Die aktuelle und geschätzte Position des Menschen, sowie die aktuelle Position des Menschen werden als Parameter übergeben, und dienen als Grundlage für die Berechnung der Trajektorie.
Die Struktur der zu berechnenden Trajektorie $s$ ist durch das Modul \lstinline{OmniDriveAlongPath} (s. Abs. \ref{bahnplanung_lowlevel_subsec}) \todoprivate{ref richtig?}\ vorgegeben:

% \begin{equation}
% 	s =
% 	\begin{pmatrix} start \\ next \\ following \end{pmatrix} =
% 	\begin{pmatrix} x_{st} & y_{st} & 0 \\ x_{nx} & y_{nx} & yaw_{nx} \\ x_{fl} & y_{fl} & 0 \end{pmatrix}
% \end{equation}
% 
% \todo[inline]{welche Variante nehmen?}

\begin{equation}
	s =	\begin{pmatrix} start \\ next \\ following \end{pmatrix}
\end{equation}
mit 
\begin{equation}
	start = 		\begin{pmatrix}	x_{start} \\ y_{start}							\end{pmatrix},
	next = 			\begin{pmatrix} x_{next} \\ y_{next} \\ yaw_{next}	\end{pmatrix},
	following =	\begin{pmatrix}	x_{following} \\ y_{following}							\end{pmatrix}
\end{equation}
Die drei Trajektorien-Punkte sind Ortsangaben mit einer $x$- bzw. $y$-Koordinate, während $following$ zusätzliche eine Komponente für $yaw$ (die ``Blickrichtung'' des Roboters) beinhaltet.\todoprivate{Skizze zu x-y-z-roll-pitch-yaw?}

Die Berechnung einer neuen Trajektorie wird durchgeführt, sobald neue Daten vom Modul \lstinline{HumanPosition} vorliegen.
Als $start$-Punkt wird dabei die aktuelle Roboterposition verwendet.
$next$ ist eine Roboter-Pose, die sich im Abstand $d_{hum}$ vor dem Menschen auf dem Vektor Roboter-Mensch befindet.
Im Abstand $d_{hum}$ vor der geschätzten zukünftigen Mensch-Position auf dem Vektor next-estimated befindet sich der dritte Punkt der Trajektorie.
Der Abstand $d_{hum}$ zum Menschen kann dabei als externer Parameter vorgegeben werden, der Wert beträgt standardmäßig $1,75m$.
In Abb. \ref{fig:trajektorie}\todoprivate{ref setzen}\ wird dies dargestellt.

\missingfigure{Vektor-Zeichnung}


\begin{itemize}
% 	\item Modul bekommt aktuelle und geschätzte Position des Menschen, sowie aktuelle Position des Roboters
% 	\item Struktur der zu berechnenden Trajektorie duch \lstinline{OmniDriveAlongPath} \todoprivate{ref zu LowLevel?}\ vorgegeben:
% 	\begin{itemize}
% 		\item Start $(x,y)$, Koordinaten
% 		\item Following $(x,y,yaw)$, Koordinaten und Orientierung \todoprivate{Skizze zu x-y-z-roll-pitch-yaw ?}
% 		\item Next $(x,y)$, Koordinaten
% 	\end{itemize}
% 	\item Berechnung einer neuen Trajektorie, sobald neue Mensch-Daten vorliegen
% 	\item Berechnung der Trajektorien-Punkte im Einzelnen:
% 	\begin{itemize}
% 	  \item Start: aktuelle Roboterposition
% 	  \item Next: Pose auf Roboter-Mensch-Vektor, die sich im Abstand $d$ vor dem Mensch befindet; yaw wird gesondert berechnet, sodass Roboter immer zu Mensch schaut
% 	  \item Following: Position auf next-estimated-Vektor, die sich $d$ vor geschätzter Mensch-Position befindet
% 	  \item erfolgt alles in absoluten Weltkoordinaten
%   \end{itemize}
%   \item Parameter $d$ für den Abstand des Roboters zum Menschen, default: $1,75m$
  \item Berechnung des yaw-Wertes: \todoprivate{Formel dazu}

% 	\item MUX für Trajektorie
% 	\begin{itemize}
% 		\item GenerateTrajectory
% 		\item LoadFromFile
% 	\end{itemize}
\end{itemize}


\todoprivate{Zeichnung mit Vektoren}



\subsubsection{Pfadspeichern und -laden}

\begin{itemize}
  \item Im \lstinline{followPerson} -Zustand: Wegpunkte $(x,y,yaw)$ werden in gewissen Abständen $d_{save}$ gespeichert
  \begin{itemize}
    \item sobald $ \overline{ s_{act} - s_{last} } > d_{save} $ \todoprivate{Fehler im Code??? Euklidischer Abstand stattdessen verwenden?}
  \end{itemize}
	\item Im \lstinline{loadPath} -Zustand: Nächster Wegpunkt wird geladen, sobald vorheriger erreicht wurde
  \begin{itemize}
		\item Im \lstinline{returnHome} -Zustand: Nächster Wegpunkt wird geladen, sobald vorheriger erreicht wurde; umgekehrte Reihenfolge
    \item sobald $ d( s_{foll} - s_{act} ) > d( s_{nxt} - s_{act} ) $, Euklidischer Abstand
  	\item \lstinline{yaw} wird unabhängig von Datei berechnet, sodass Roboter immer Menschen anschaut
  \end{itemize}
  \item Eigene Klasse zum File-Handling: \lstinline{fileUtil}
  \begin{itemize}
    \item Mehrere dump-files möglich
    \item Auslesen der Wegpunkte in normaler und umgekehrter Reihenfolge
  \end{itemize}
\end{itemize}



% 
% % -----------------------------------------------------------------------------
% %													Transformationen
% \subsection{Transformationen}
% \label{transformationen_subsec}
% 
% \authorsection{\editorjulian}
% \todo[inline]{Julian: Sollen wir (also du ;-) ) noch so einen Abschnitt machen?}


% -----------------------------------------------------------------------------
%													Kollisionsvermeidung
\subsection{Kollisionsvermeidung}
\label{kollisionsvermeidung_subsec}

\subsubsection{Modul}
\authorsection{\editortobias}
\todo[inline]{Tobi: Schreiben}

\begin{itemize}
 
  \item Kollisionsvermeidungsmodul (\lstinline{ObstacleAvoidance}) überprüft die geplante Trajektorie auf Hindernisse, und plant ggf. um Hindernisse
  \item Ausgaben von \lstinline{GenerateTrajectory} und \lstinline{LoadFromFile} werden geMUXt
  \begin{itemize}
    \item je nach Modus (\lstinline{followPerson} bzw. \lstinline{returnHome} / \lstinline{followPath}) schaltet MUX andere Eingänge durch
  	\item Die Ausgabe des MUX wird dann widerum komplett durch das Kollisionsvermeidungsmodul (\lstinline{ObstacleAvoidance}) geschleift
    \item Dadurch wird jede Trajektorie auf Hindernisse überprüft, keine kommt \glqq daran vorbei\grqq
  \end{itemize}
  \item Innerhalb des Moduls wird jede Trajektorie auf Hindernisse überprüft
	\begin{itemize}
	  \item die (Bild-)Punkte zwischen start, next und following werden berechnet (Transformation der Welt- in Bildkoordinaten)
	  \item Überprüfung für jeden einzelnen Punkt, ob er innerhalb eines Hindernisses liegt
	  \item Verwendung der \lstinline{ObstacleMap} \todoprivate{ref zu Lokalisierung}
	\end{itemize}
  \item Ohne Hindernis werden Trajektorien-Punkte nur durchgeschleift
  \item Wenn Hindernis erkannt, dann wird eine neue Trajektorie ermittelt
  \begin{itemize}
    \item Verwendung des A*-Algorithmus
    \item Ermittlung des Start- und Endpunktes für A*
    \item ggf. gesonderte Behandlung für Endpunkt (\lstinline{choosePointAfterObstacle})
    \item Aufruf A*
    \item Transformation der nuen Trajektorie IK $\rightarrow$ WK
  \end{itemize}

\end{itemize}


\subsubsection{Algorithmen}
\authorsection{\editorjulian}
\todo[inline]{Julian: Schreiben}

Um einen kollisionsfreien Pfad zu erstellen kommt eine Bahnplanungsmethode zum Einsatz, die den A*-Algorithmus verwendet. Um den zeitlichen Aufwand einer eigenen Implementierung zu vermeiden wurde auf eine bereits existierende und frei verfügbare Implementierung von Justin Heyes-Jones zurückgegriffen.\todoprivate{Website und ref angeben}


\begin{itemize}
		\item Anpassungen an bestehender Implementierung: Bitmap karte, Heuristik
		\item Erweiterung durch Gravitationsfeld
\end{itemize}

% -----------------------------------------------------------------------------
%													Simulation
\subsection{Simulation}
\label{simulation_subsec}
\authorsection{\editoroier}

Für die Simulation war schon eine Simulationsumgebung vorhanden, Aufgabe war eine grafische Oberfläche mit Mcagui zu basteln.
Warum ist die Simulation nützlich oder notwendig? In erster Linie konnte man am Anfang an dem realen Roboter nicht testen, weil der Roboter noch gebaut wurde oder zum Beispiel nicht alle von uns programmierte Komponente, wie bspw. die Lokalisierung, funktionsfähig waren. Dank der Simulation wurde das testen von einzelne Funktionalitäten und die Darstellung des Roboterzustands möglich wie man auf der Abbildung \ref{fig:mcagui} sehen kann. Zum Beispiel wurde die Simulation sehr oft benutzt um schnell Änderungen im Bahnplanungsalgorithmus zu testen.

\begin{figure}[h]
\centering
\includegraphics[scale=0.4]{graphics/mcagui_screenshot.png}		
\caption{\label{fig:mcagui} Die entwickelte Simulationsoberfläche
}
\end{figure}

Ein Problem am Anfang war, wie man den zu verfolgenden Mensch simulieren sollte. Um dies zu lösen, benutzt man einfach Mausklicks auf der Karte.

% #############################################################################
%															Hinderniskarte und Lokalisierung
% #############################################################################
\chapter{Hinderniskarte und Lokalisierung}
\label{lokalisierung_cha}
% ********************************************************************************
% 										Aufgabenstellung
% ********************************************************************************
\section{Aufgabenstellung}
\label{lokalisierung_aufgabenstellung_sec}
\authorsection{\editorandreas}
% -----------------------------------------------------------------------------
%													Lokalisierung
\subsection{Lokalisierung}
Die Aufgabe der Lokalisierung, im Kontext der Robotik, ist die Bestimmung der
Roboter Pose, welche die Position und Orientierung des Roboters in seiner Umgebung vollständig beschreibt.
 Um Bahnplanung und Steuerung zu ermöglichen muss die Pose des Roboters bekannt sein,
 damit sich dieser gefahrlos fortbewegen und sicher mit der Umwelt interagieren kann.
 Für einen Menschen ist die Lokalisierung selbstverständlich, für einen Roboter
 stellt sie allerdings eine Herausforderung dar.
 Aufgrund einer unbekannter Startposition oder von Messungenauigkeiten während
 der Fortbewegung, kann die Pose eines Roboters nur in den seltensten Fällen
 genau bestimmt werden.
 Von lokaler Lokalisierung wird gesprochen wenn die aktuelle Pose des Roboters
 in seiner Umwelt bekannt ist und bei einer Positionsänderung des Roboters diese
 fortlaufend aktualisiert wird, um sie auf dem neusten Stand zu halten. Hier
 werden Schätzungen der Bewegung aufgrund der Odometrie Sensorik mit Hilfe von weiteren Daten zur Posenbestimmung korrigiert. Dieses Vorgehen soll eine Aufsummierung von Fehlern bei der
 Schätzung der Bewegung durch die Odometrie verhindern. Im Gegensatz dazu ist
 bei der globalen Lokalisierung die aktuelle Pose des Roboters in seiner Umwelt
 nicht bekannt und muss zum Beispiel anhand von Kartenmaterial geschätzt werden.
 Man sieht leicht ein, dass bei dieser Art der Lokalisierung die Initiale Schätzung einen beliebig großen Fehler aufweisen kann.
 Nach dieser Einführung sollte klar geworden sein das es notwendig ist eine Lokalisierung durchzuführen um es unserem Roboter zu ermöglichen sich fortzubewegen.
% -----------------------------------------------------------------------------
%													Hinderniskarte
\subsection{Hinderniskarte}
Abgesehen von der Lokalisierung gibt es noch einen weiteren Grund warum es für einen mobilen Roboter unumgänglich ist,
 dass eine interne Repräsentation seiner Umgebung existiert. Karten, die zur Lokalisierung eingesetzt werden,
 enthalten meist nicht alle statischen Hindernisse, da diese Karten zu einem gewissen Zeitpunkt erstellt wurden und
 in der Zwischenzeit wahrscheinlich neue Objekte hinzugekommen oder bereits Vorhandene versetzt worden sind.
 Außerdem enthalten solche Karten offensichtlich keine Hindernisse die sich bewegen können, wie beispielsweise Menschen.
 Ein Hindernis ist ganz allgemein ein Bereich der von dem Roboter nicht befahren werden kann. Eine Darstellung der Hindernisse
 relativ zur Roboterpose wird als Hinderniskarte bezeichnet. Eine solche Karte
 kann zum Beispiel mit Hilfe von Laser-- oder Kameradaten erstellt werden.
 Eine Kombination des Materials mehrerer Sensoren ist normalerweise das Sinnvollste,
 da Laser beispielsweise keine negativen Hindernisse erkennen können.
 Als negative Hindernisse werden Bereiche bezeichnet die im Vergleich zur
 gegenwärtigen Position eine geringere Höhe aufweisen.
 Ein Beispiel für ein solches Hindernis ist eine Treppe in ein tiefer gelegenes
 Stockwerk. Das Erstellen einer Hinderniskarte ist aus den oben genannten
 Gründen essentiell um eine Bahnplanung für einen mobilen Roboter
 durchführen zu können.
% -----------------------------------------------------------------------------
%													Virtual protective Field
\subsection{Virtual protective Field}
Die Aufgabe der Bahnplanung besteht darin auf Grundlage der Lokalisierung und der Hinderniskarte einen Weg zu einer
 bestimmten Position zu finden. Danach wird dieser Weg vom Roboter abgefahren. Allerdings kann es während dem Abfahren der gefundenen
 Route Aufgrund eines Hindernisses, dass sich in den Weg des Roboters bewegt bevor erneut geplant werden kann oder
 schlicht Aufgrund eines Fehlers in der Planung dazu kommen, dass der Roboter mit einem Hindernis kollidiert.
 Daher wird eine Möglichkeit benötigt auf einer tieferen Ebene wie der Bahnplanung das unbeabsichtigte kollidieren mit einem Hindernis
 zu verhindern. An dieser Stelle kommt das sogenannte „Virtual protective Field“ ins Spiel, von welchem ein Bereich um den Roboter auf
 Hindernisse überprüft und beim Auffinden eines Solchen ein Signal erzeugt wird, damit der Roboter gestoppt werden kann bevor er
 auf das Hindernis trifft. Dieser Bereich muss mit der Geschwindigkeit des
 Roboters wachsen bzw. schrumpfen so das sichergestellt ist, dass auf jeden
 Fall eine Kollision verhindert werden kann.
% ********************************************************************************
% 										Grundlagen Lokalisierung
% ********************************************************************************
\section{Grundlagen Lokalisierung}
\label{lokalisierung_grundlagen_sec}
\authorsection{\editoranne}
\begin{flushright}\textit{Einige Abschnitte in Anlehnung an das oder aus dem Projektwiki von Dirk~Toewe}\end{flushright}
% -----------------------------------------------------------------------------
%													Lokalisierung
\subsection{Lokalisierung}
Um die im vorangegangenen Abschnitt vorgestellten Aufgaben und Ziele umzusetzen, gibt es verschiedene Möglichkeiten. Wir gehen nun etwas detailierter auf grundlegende Verfahren ein, die für die Durchführung einer Positionsbestimmung von Nutzen sind.
 
\subsubsection{Odometrie}
Die erste vorgestellte Methode zur Positionsbestimmung ist die Odometrie. Dieser aus dem griechischen kommende Begriff steht sinngemäß für ''Wegmessung'' und beschreibt im Wesentlichen schon die Hauptaufgabe. Die Odometrie umfasst die technischen Einrichtungen, die während des Betriebes des Roboters die zurückgelegte Wegstrecke ausmessen.

Ein sehr einfacher und dennoch effektiver Ansatz zur Wegbestimmung ist die Verwendung von Schrittzählern. Bewegt sich der Roboter mittels Beinen fort, wird die Anzahl der gemachten Schritte gezählt. Werden dagegen Rädern  zur Fortbewegung genutzt, können schwarz--weisse Musterscheiben angebracht werden. Diese werden mittels einem oder mehrerer Sensoren abgetastet. Eine einfache, zumeist lineare, Abbildung ermittelt aus der Anzahl der schwarzen und weissen Flächen die zurückgelegte Wegstrecke. Es kann anschließend relativ zur Ausgangpose eine neue Roboterpose angegeben werden.

\begin{figure}[h]
\center
\includegraphics[scale=0.5]{graphics/odo_einbau.jpg}
\caption{\label{fig:Odo_Einbau} Beispiel einer eingebauten Musterscheibe an einer Radachse\citep{odom}}
\end{figure}

Die Vorteile liegen bei diesem Ansatz in der Einfachheit der Konstruktion. Der benötigte Aufwand zur Installation, bzw. Implementation ist nicht sehr hoch und auch nach dem Einbau ist der Rechenaufwand vernachlässigbar gering, da keine komplizierten Funktionsberechnungen nötig sind. Die Nachteile der Methode beruhen vor allem auf den Fehlern, die mit dieser Methode entstehen. Ist der Boden uneben oder besitzt zu viel Schlupf sind die Radumdrehungen nicht mehr konsistent zur zurückgelegten Strecke. Auch die Annahmen über die Rad-- und Fahrgestellgeometrie können Ungenauigkeiten enthalten. So ist zum Beispiel ein Rad nicht gänzlich rund oder die Radabstände wurden ungenau gemessen oder haben sich durch Verschleiß und Benutzung verändert. Zu diesen Nachteilen, die an sich schon eine fehlerhafte Wegbestimmung nach sich ziehen, kommt noch erschwerend hinzu, dass die Odometrie keine rückgekoppelte Methode ist, sich eingeschlichene Fehler selbstverständlich mit dem Betrag der zurückgelegten Wegstrecke aufsummieren und dies auch noch abhängig von der aktuellen Situation in unterschiedlichem Maße. Der Fehler kann also nicht fest mit ein-- oder gar gänzlich herausgerechnet werden. Des Weiteren kann nur eine relativ zu einer Ausgangspose berechnete neue Roboterpose angegeben werden. Diese Methode eignet sich daher nur für eine lokale, nicht jedoch für eine globale Lokalisierung. Die Ursprungsposition muss dem Roboter von aussen mitgeteilt werden.

\subsubsection{Kartenbasierte Lokalisierung}

Ist eine Karte der Umgebung gegeben, in der sich der Roboter aufhält, so kann mittels geeigneter Sensordaten die Roboterpose geschätzt, bzw. eindeutig bestimmt werden. Eine Möglichkeit besteht darin, Laser an dem Roboter anzubringen und mittels Abstandsdaten Gegenstände, Wände und Flächen zu erkennen. Dazu wird aus einer Punktwolke eine Menge von geometrischen Figuren, wie zum Beispiel Kanten extrahiert. Diese werden mit den Kartenmaterial abgeglichen und man erhält im schlechten Fall mehrere mögliche Positionen, im optimalen Fall eine eindeutige Position an der sich der Roboter aktuell befindet.

Hierbei wird immer die aktuelle Roboterpose berechnet ohne auf vorhergehende Roboterposen angewiesen zu sein. Dies ist vorteilhaft, da sich im Laufe des Betriebes keine Fehler akkumulieren können und man nicht einmal eine Ausgangspose benötigt. Allerdings ist der Roboter auf möglichst detailgetreues Kartenmaterial angewiesen, welches selbstverständlich aktuell gehalten werden muss. In einer Umgebung in der keine markanten Stellen zu finden sind, wird die Zahl der möglichen Roboterposen eventuell unangenehm hoch sein. Dazu kommen die Nachteile der Technik, die verwendet wird um Informationen über die Umwelt zu erhalten. Bei einem Laser muss darauf geachtet werden, dass er freie Sicht hat und gerade am Roboter angebracht ist. Zusätzlich kann der Laser fehlerhafte Daten zurückliefern, wenn er auf spiegelnde, glatte Flächen trifft, durch eine Glaswand hindurchsieht oder von sich bewegenden Gegenständen umgeben ist.  

\subsubsection{Landmarkbasierte Lokalisierung}

In einer bekannten Umgebung können Landmarks (\gls{qr}--Codes, Bodenlinien,
 \gls{gps}--Satelliten) angebracht werden. Diese können durch Sensoren
 am Roboter wiedererkannt werden und ermöglichen die Lokalisierung
 z.B. mittels Triangulierung.

\begin{itemize}
  \item Vorteile:
  \begin{itemize}
    \item (durchschnittlicher) Fehler niedrig
    \item Zuverlässig (max. Fehler niedrig)
    \item Keine Fehlerakkumulation über Zeitverlauf
  \end{itemize}
  \item Nachteile:
  \begin{itemize}
    \item Landmarks notwendig
    \item Landmarks müssen von überall sichtbar sein
    \item Kartenmaterial u.U. notwendig
   \end{itemize}
\end{itemize}

Um die Nachteile der verschiedenen Lokalisierungsmethoden zu
 verringern werden in der Regel mindestens zwei Methoden miteinander
 kombiniert. Aufgrund der einfachen Implementierung und der geringen
 Störanfälligkeit ist meistens die Odometrie in Kombination mit
 anderen Systemen vorzufinden. Bei Verwendung mehrerer Methoden sind
 dementsprechend viele fehlerbehaftete Annahmen der aktuellen
 Roboterpose gegeben. Im Folgenden wollen wir zwei unterschiedliche
 Verfahren vorstellen, welche aus diesen Annahmen bestmöglich die
 aktuelle Roboterpose bestimmen.
 
\subsubsection{Kalman--Filter}

Das Kalman--Filter ist eine der wichtigsten Entdeckungen des letzten Jahrhunderts, entwickelt von dem Ingenieur Rudolf E. Kalman.
Sind eine Reihe von Messwerten in Form von möglichen Roboterposen gegeben, die von unterschiedlichen Messverfahren herrühren und mit einem Rauschen behaftet sind, so schätzt das Kalman--Filter eine aktuelle Roboterpose.
Die Grundidee sieht man an der Ausgangsgleichung \[ \hat{p}_k = G_k m_k + (1-G_k) \hat{p}_{k-1}, \] wobei \(\hat{p}_k\) der aktuelle und  \(\hat{p}_{k-1}\) der vorherige Schätzwert sei. \(G_k\) das sogenannte Kalman--Gain und \(m_k\) der aktuelle Messwert, der allerdings verrauscht ist und daher mit einem Fehler behaftet. Das Kalman--Gain gewichtet die bisherigen Messwerte in Bezug auf den akuell gemessenen Wert und berechnet daraus den aktuellen Schätzwert \(\hat{p}_k\). Das Ziel ist einen möglichst guten Schätzwert \(\hat{p}_k\) zu finden, der den Wert \(m_k\) am besten beschreibt.
In den Iterationsschritten wechseln sich nun Prädiktion und Korrektion ab: Aus den letzten \(k-1\) Messungen wird ein vorraussichtlicher Wert \(\hat{p}_k\) vorhergesagt (''Prädiktion''). Nachdem die \(k\)te Messung stattgefunden hat, man also eine Information mehr dazu gewonnen hat, wird der Messwert zum Zeitpunkt \(k\) verbesserert (''Korrektur''). Dabei kann in jedem Schritt aus den gemachten Beobachtungen das Kalman--Gain neu angepasst werden.

\subsubsection{Partikel--Filter}

Als Alternative zum Kalman--Filter existiert die Sequentielle Monte--Carlo--Methode. Auch bekannt als Partikel--Filter. Mit Hilfe dieser Methode können ebenfalls kontinuierlich Ort und Geschwindigkeit, insbesondere eines autonomen Roboters, ermittelt werden, obwohl nur ungenaue Messwerte vorliegen.
Ist das Prozessmodell nicht linear oder sind die äußeren Einflüsse, die zu fehlerhaften Messwerten führen nicht normalverteilt, kann auch nicht von einer Normalverteilung des zu schätzendes Zustandes ausgegangen werden. Eine Normalverteilung hat die Eigenschaft nur ein Maxima aufzuweisen und somit nur einen möglichen besten Zustand zurück zu liefern. Liegt diese nicht vor, muss angenommen werden, dass die Wahrscheinlichkeitsdichte mehrere Maxima aufweist, das heisst, das mehrere mögliche Zustände in Frage kommen, als die Besten und Wahrscheinlichsten zum gegenwärtigen Zeitpunkt zu gelten.
Diese unbekannte, aber gerade aktuelle Warscheinlichkeitsdichte auf dem Zustandsraum soll nun geschätzt werden, um daraus den wahrscheinlichsten Zustand zu schätzen. Dazu wird eine Menge von Partikeln erzeugt. Jeder Partikel ist dabei ein Paar aus Gewicht und einem Punkt im Zustandsraum. Die Menge der Partikel repräsentiert die Wahrscheinlichkeitsdichte.
Jedem Partikel wird dann eine mögliche Lösungskurve zugeordnet und je besser diese zu den tatsächlichen Messwerten passt, desto mehr wird das Gewicht des Partikels angepasst. Im Laufe der Messungen verbessert sich somit die Schätzung der Wahrscheinlichkeitsdichte, welches durch Methoden der nichtparametrischen Dichteschätzung erfolgt.\citep{rob3}

 
% ********************************************************************************
% 										Umsetzung
% ********************************************************************************
\section{Umsetzung}
\label{lokalisierung_umsetzung_sec}
\authorsection{\editorandreas}
\begin{flushright}\textit{Einige Abschnitte aus dem Projektwiki übernommen deren
Autor: Dirk~Toewe}\end{flushright}
% -----------------------------------------------------------------------------
%													Lokalisierung
\subsection{Lokalisierung} 
In diesem Abschnitt wollen wir uns mit der eigentlichen Umsetzung der
 Lokalisierung beschäftigen. Wie sie wahrscheinlich schon aus Kapitel
 \ref{grundlagen_cha} wissen, verfügen die beiden Roboter \gls{hollie} und
 \lstinline{Odete} über sehr ähnliche Sensoren, die zur Lokalisierung verwendet
 werden können. Der Roboter \lstinline{Odete} weist 2 Winkelencoder und einen 270 Grad Sick
 Laserscanner (Scheibenförmiger Scan) auf.
 \gls{hollie} verfügt über 2 Winkelencoder und zwei
 270 Grad Sick Laserscanner.
 Offensichtlich musste, um ein sicheres Fortbewegen des Roboters zu
 gewährleisten, für das \lstinline{Odete}--Projekt ebenfalls eine Lokalisierung implementiert werden.
 Im Folgenden wollen wir kurz auf die grundlegende Idee für den Aufbau der
 Gruppen und Module bezüglich der Lokalisierung eingehen.
 Die verwendete \lstinline{MCA2}--Struktur der \lstinline{Odete} basiert auf
 einem bestimmten Paradigma.
 Ziel ist es eine hohe Flexibilität zu erreichen. Sensoren und Plattformen sollen
 beliebig kombinierbar sein. Wir entschieden uns daher dazu diese Struktur auch
 beim \lstinline{SegwayOmni}--Projekt weiterzuführen. Dadurch konnten alte
 Programmstrukturen übernommen und/oder adaptiert und außerdem dieselbe
 Flexibilität wie im \lstinline{Odete}--Projekt erreicht werden. Daher gelten die folgenden Erläuterungen
 zur Lokalisierung auch fast uneingeschränkt für beide Robotersysteme.
 Zur Positionsschätzung wird bei beiden Systemen eine Kombination aus
 kartenbasierter Lokalisierung und Odometrie verwendet. Die Schätzwerte aus beiden Messungen werden mit einem Kalman--Filter zu einer
 Synthese geschätzt. Dem Kalman--Filter ist eine Queue (\gls{fifo}) vorgeschaltet, in der jede Schätzung
 mit einem Timestamp registriert wird. So können die Messwerte in zeitlich korrekter Reihenfolge
 eingerechnet werden, selbst dann, wenn sie nach der Verarbeitung in einer anderen Reihenfolge
 im Kalman--Filter registriert werden. Die für die Odometrie benötigte Winkelschrittmessung der Räder wird direkt
 von der Segway--Antriebsplattform geliefert.
 Es wird angenommen, dass der Fehler linear ansteigt mit der durch den Reifen zurückgelegten Strecke.
 Falls eine gute Schätzung aus der kartenbasierten Lokalisierung in den Kalman--Filter übergeben wird,
 wird die Pose korrigiert und der Odometriefehler wieder gesenkt.
 Nach den Erläuterungen dazu wie die Lokalisierung durchgeführt wird wollen wir
 jetzt die beteiligten Gruppen und Module etwas näher betrachten. 

\begin{description}
\item[HAL(Hardware abstraction Layer)] Die \lstinline{Hal}--Gruppe ist für
 plattformabhängige Hardware zuständig.
 Mit der Plattform sind jene Hardwarekomponenten gemeint, die in der Regel nicht getauscht werden 
 (Rechner, Omnidirektionaler Antrieb, Akummulatoren).
 In \lstinline{MCA2} umfasst die \lstinline{Hal} die Bewegungsplattformen, die dazugehörigen
 Winkelencoder und die Posenerfassung (nicht Posenberechnung) des Roboters. Die \lstinline{Hal} ist eine Abstraction Layer in dem Sinne, dass zwischen einer Simulierten
 Hardware und der tatsächlichen Hardware gewechselt werden kann, ohne das andere
 \lstinline{MCA2}--Gruppen angepasst werden müssen.
 Prinzipiell können sogar die Plattformen (Omnidirektional <-> Differentiell) ausgetauscht werden. Die Odometrie hängt
 direkt mit der Antriebsart zusammen und findet deswegen in der \lstinline{Hal} statt.
 Auch die Kalman--Filterung findet in dieser Gruppe im
 \lstinline{AbsoluteCoordinatesMecanumKLF}--Modul statt (KLF = KalmanFilter). Für
 eine Darstellung der Gruppe siehe Abbildung \ref{fig:segwayHal}.
\item[SAL(Scanner abstraction Layer)] Die \lstinline{Sal}--Gruppe umfasst alle ``angebauten''
Scanner, nicht aber in Plattform integrierte Scanner (Winkelencoder).
 Im Sinne der Abstraction Layer können auch hier Hardware-- und
 Simulationsversion der \lstinline{Sal} ausgewechselt werden, siehe Abbildung
 \ref{fig:SegwayOmniSal}
\end{description}
\begin{figure}[h]
\center
\includegraphics[scale=0.7]{graphics/SegwayOmniSal.jpg}
\caption{\label{fig:SegwayOmniSal} SegwayOmniSal}
\end{figure}

Diese Beiden Gruppen realisieren also unter anderem die erste benötigte
 Abstraktion der Sensordaten sowie eine Simulation für den Testfall. Zu beginn
 des Projektes existierte bereits eine \lstinline{Hal} Gruppe im
 \lstinline{SegwayOmni}--Projekt, allerdings war darin noch keine kartenbasierte
 Positionskorrektur durch den Kalman--Filter realisiert. Diese Funktionalität und
 die \lstinline{Sal}--Gruppe wurden einige Zeit später durch einen unserer Betreuer,
 Jan Oberländer implementiert. Da wir Anfangs nicht wussten
 welche Funktionalitäten von uns implementiert werden mussten hatte unserer Gruppe einige Schwierigkeiten
 einen Einstieg zu finden. Dies äußerte sich vor allem darin, dass wir zu Beginn des Praktikums viel Zeit damit
 verbrachten zu versuchen uns in alle Gruppen und Module einzuarbeiten.

Nach der Feststellung der benötigten Anpassungen an die vorhandenen Gruppen und
 Module des \lstinline{Odete}--Projekts wurden diese adaptiert. Außerdem wurden
 Erweiterungen integriert damit eine beliebige Menge von Scannern genutzt werden
 kann.
 Eine Parametrisierung über einen Multiplexer und ein XML--Config--File
 ermöglichen eine Anpassung auf eine neue Sensor--Anordnung. Abschließend erfolgt
 nun noch eine detailliertere Erläuterung der für die kartenbasierte
 Lokalisierung nötigen Schritte und der involvierten Gruppen und Module. Die kartenbasierte Lokaliserung
 findet in den \lstinline{MCA2}--Gruppen \lstinline{LlSP} und \lstinline{HLSP}
 statt. Die Laserscanner--Daten dazu werden wie schon erörtert aus der
 \lstinline{SAL}--Gruppe bezogen.
 Für eine Darstellung der beiden Gruppen im \lstinline{Mcabrowser} siehe
 Abbildung \ref{fig:Llsp} und \ref{fig:Hlsp}.
 
\begin{figure}[h]
\center
\includegraphics[scale=0.7]{graphics/Llsp.jpg}
\caption{\label{fig:Llsp}Llsp -- Low Level Scanner Processing}
\end{figure}

\begin{figure}[h]
\center
\includegraphics[scale=0.7]{graphics/Hlsp.jpg}
\caption{\label{fig:Hlsp}Hlsp -- High Level Scanner Processing}
\end{figure}
 
\begin{description}
\item[LLSP(Low Level Scanner Processing)] Die \lstinline{LLSP} dient der
Verarbeitung der Scanner Rohdaten. Im Gegensatz zur \lstinline{HLSP} macht die
\lstinline{LLSP} keine logische Auswertung der Scannerdaten.
 Es werden lediglich grundlegende Operationen auf die einzelnen Messwerte angewandt
 (Koordinatentransformationen, Löschen ungültiger Werte, ...).
\item[HLSP(High Level Scanner Processing)] In der \lstinline{HLSP} werden
logische Operationen auf die\\ Messwerte angewandt (Kantenextraktion,
Posenkorrektur--Schätzung).
\begin{itemize}
  \item Der \lstinline{ScanStrategyMultiplexer} dient als zentrale Schaltstelle
  für die Controllerinputs der \lstinline{HLSP}--Module.
  Über einen einzelnen Controllerinput des \lstinline{ScanTrategyMultiplexer}
  kann ein Setup geladen werden, welches die Controllerinputs der einzelnen
 \lstinline{HLSP}--Module setzt.
\end{itemize}
\end{description}

Die kartenbasierte Posen--Schätzung wird in den folgenden Schritten ermittelt:

\begin{enumerate}
  \item Scheibenförmiger Laser--Scan der Umgebung parallel zur Fahrbahn\\
    	Gruppe: \lstinline{SAL}\\
    	Ergebnis: Winkel--Abstand--Paare (Polarkoordinaten)\\
  \item Koordinaten--Transformation: polare Sensorkoordinaten -> kartesische
  		Sensorkoordinaten\\
  		Gruppe: \lstinline{LlSP}\\
    	Modul: \lstinline{ScanToRobotCS}\\
    	Ergebnis: XY--Punkte (kartesische Sensorkoordinaten)\\
  \item Koordinaten--Transformation: kartesische Sensorkoordinaten -> kartesische
  		Roboterkoordinaten\\ 
  		Gruppe: \lstinline{LlSP}\\
  		Modul: \lstinline{MergeScan}\\
  		Ergebnis: XY--Punkte (kartesische Roboterkoordinaten)\\
  \item Löschen von Punkten mit zu hohem Abstand oder schlechter Auflösung\\
    	Gruppe: \lstinline{LlSP}\\
    	Modul: \lstinline{CutScan}\\
    	Ergebnis: XY--Punkte mit gültigem Abstand\\
  \item Extraktion/Schätzung von Kanten aus der Punktemenge\\
    	Gruppe: \lstinline{HlSP}\\
    	Modul: \lstinline{ExtractEdges}\\
    	Ergebnis: Liniensegmente\\
  \item Vergleich der Kanten mit der Karte und Schätzung einer Pose samt
  		Fehlerkovarianz\\ 
  		Gruppe: \lstinline{HlSP}\\
    	Modul: \lstinline{PoseCorrection}\\
    	Ergebnis: X,Y,Gierwinkel--Schätzung + Fehlerkovarianz\\
\end{enumerate}

% -----------------------------------------------------------------------------
%													Hinderniskarte
\subsection{Hinderniskarte}
\label{hinderniskarte_subsec}
\authorsection{\editorandreas}
 Im Gegensatz zur Lokalisierung gab es für die von uns gewünschte Lösung dieser
 Aufgabenstellung noch keine Gruppen oder Module von älteren Projekten. Aus
 diesem Grund war es notwendig im Vorfeld intensivere Überlegungen über die möglichen
 Umsetzungen und Repräsentationen anzustellen.
 Nach einigen Diskussionen entschieden wir uns für die Nutzung von
 Bildverarbeitungs"=Modulen aus der \lstinline{SenseIVT}--Bibliothek
 und einer Repräsentation der Hinderniskarte durch ein
 Schwarzweiß--Bitmap. Der Grund für diese Wahl lag vor
 allem an der guten Visualisierung, die dadurch erhalten werden kann, außerdem wird eine einfache
 Anpassung und die Verwendung von Standard--Filtern ermöglicht.
\begin{figure}[h]
\center
\includegraphics[scale=0.7]{graphics/hinderniskarte.jpg}
\caption{\label{fig:hinderniskarte} Hindernisskarte als Schwarzweiß--Bitmap}
\end{figure}
 In Abbildung \ref{fig:hinderniskarte} ist ein Beispiel eines solches Bitmaps
 zu sehen. Mit der Semantik, dass weiße Areale Hindernisse darstellen und schwarze Bereiche entweder befahrbar
 oder von dieser Position aus nicht einsehbar sind. Diese von uns gewählte Abbildung der Werte
 auf die entsprechenden Bedeutungen ist allerdings beliebig und hat keinen besonderen Grund.
 Die Roboterposition ist der Mittelpunkt der Karte. Bevor Hindernisse in eine solche Karte eingezeichnet werden
 können müssen sie offensichtlich zuvor durch Sensoren detektiert worden sein. Es lag nahe die Laserdaten,
 welche auch zur Lokalisierung eingesetzt und bereits in geeigneter Form vorhanden waren zu nutzen.
 Da ein Roboter in der Realität eine größere Fläche als ein Punkt aufweist entschieden wir uns bereits an dieser
 Stelle zu einer Expansion der Hindernisse auf Robotergröße, damit solche Berechnungen an anderer Stelle vermieden
 werden können.
\begin{figure}[h]
\center
\includegraphics[scale=0.8]{graphics/ObstacleExtractor.jpg}
\caption{\label{fig:obstacleExtractor} Gruppe \lstinline{ObstacleExtractor}}
\end{figure}
 In Abbildung \ref{fig:obstacleExtractor} ist die die Gruppe
 (\lstinline{ObstacleExtractor}), durch welche die gewünschte Funktionalität implementiert wurde zu sehen.
 Im Folgenden wollen wir auch hier die einzelnen Module kurz vorstellen:

\begin{description}
\item[SegwayOmniOESMultiplexer] Dieses Module ermöglicht es die anderen Module der Gruppe zu steuern.
 Hier kann durch einen Controllerinput die gesamte Gruppe aktiviert oder deaktiviert werden.
 Außerdem kann die Größe der Erode-- und Dilate--Filtermasken manipuliert werden.
\item[UpdateManager] Dieses Module ist ein Teil der
\lstinline{SenseIvt}--Bibliothek und dient dazu die Aktualisierung der Quelldaten
zu beeinflussen.
\item[SenseBitmapSource] Auch dieses Modul ist ein abgewandelter Teil der
\lstinline{SenseIvt}--Bibliothek, es dient in unserem Fall dazu ein, in einem
 XML--Config--File spezifiziertes schwarzes Bitmap zu laden.
 Dafür stehen unterschiedliche Größen zur Verfügung.
\item[SenseIvtScanToBitmap] Hier werden die Hindernisse aus den Laserdaten extrahiert und auf Robotergröße expandiert.
\item[SenseIvtDilate] Ein Dilate--Filter.
\item[SenseIvtErode] Ein Erode--Filter. Zusammen mit dem vorigen Dilate--Filter
wird hier eine Schlie\-ßen--Operation ausgeführt.
\item[ObstacleMap] Diese Modul dient dazu mit Hilfe eines Plugins des
\lstinline{Mcabrowser} die resultierende Hinderniskarte betrachten zu können.
\end{description}

% -----------------------------------------------------------------------------
%													Virtual protective Field
\subsection{Virtual protective Field}
\label{sec:vpf}
Ein Virtual protective Field, was die letzte von uns zu lösenden Aufgabe
darstellte war bereits im \lstinline{Odete}--Projekt realisiert,
 allerdings war dieses lediglich für einen Differentialantrieb implementiert.
 Nach einer Begutachtung des bestehenden Moduls entschieden wir uns dafür
 dieses nicht für omnidirektionale Bewegung anzupassen sondern eine neue Gruppe mit der gewünschten Funktionalität
 zu erstellen, da uns dies effizienter erschien. Die Wahl der für die Umsetzung benötigten Primitiven viel, aufgrund
 der von uns gewonnen Erfahrungen mit den Bildverarbeitungs--Modulen und der Möglichkeit die bereits bestehende
 Hinderniskarte auszunutzen auch hier auf die \lstinline{SenseIVT}--Bibliothek.
 Als Repräsentation wurde wieder ein Schwarzweiß--Bitmap gewählt. Auf Anfrage der
 Bahnplanungs und Steuerungs Gruppe implementierten wir 3 unterschiedliche
 Virtual protective Fields siehe Abbildung \ref{fig:virtualprotectivefields}.
 \begin{figure}[h]
\center
\includegraphics[scale=0.5]{graphics/virtualprotectivefields.jpg}
\caption{\label{fig:virtualprotectivefields} Darstellung der unterschiedlichen Protective
Fields}
\end{figure}
 Die Varianten \emph{Sideward} sowie \emph{Forward} sollen es \gls{hollie}
 ermöglichen eine Bewegung in zumindest eine Richtung durchzuführen wenn die Kombination zu einer Kollision führen würde.
 Die Virtual protective Fields unterliegen einer dynamischen Ausdehnung zum
 Quadrat der Geschwindigkeit um ein erfolgreiches Anhalten vor einer Kollision zu gewährleisten. Aufgrund der Nutzung von Bildverarbeitungs--Modulen
 konnte die Erkennung einer möglichen Kollision durch eine AND--Operation auf der Hinderniskarte und dem entsprechenden
 Virtual protective Field realisiert werden. Sollten nach dieser Operation noch
 weiße Pixel im resultierenden Bitmap vorhanden sein gilt dies als erkannte Kollision. Aufgrund von Fehlern der Laserdaten wird verlangt, dass mehr als ein
 weißer Pixel vorhanden ist, um das Verfahren etwas robuster zu machen.
\begin{figure}[h]
\center
\includegraphics[scale=0.29]{graphics/VirtualProtectiveBitmap.jpg}
\caption{\label{fig:virtualprotectivebitmap} Gruppe \lstinline{VirtualProtectiveBitmap}}
\end{figure}
 Wie bereits in den beiden vorigen Abschnitten  ist auch hier in Abbildung
 \ref{fig:virtualprotectivebitmap} die beschriebene Gruppe dargestellt und es
 folgt eine nähere Betrachtung der Module.
 Wobei wir allerdings nicht erneut näher auf die Module, die bereits zuvor
 vorgestellt wurden eingehen werden.

\begin{description}
\item[SegwayOmniVPBMultiplexer] Ähnlich wie das Modul
\lstinline{SegwayOmniOESMultiplexer} aus dem vorigen Abschnitt werden die
 übrigen Module der Gruppe hier gesteuert und können beispielsweise aktiviert
 beziehungsweise deaktiviert werden.
\item[SenseIvtProtectiveFieldSource] Dieses Modul dient in der jeweiligen
Variation dazu das entsprechende Virtual protective Field in ein schwarzes
Bitmap einzuzeichnen. Wie bereits besprochen ist die Größe des Feldes abhängig
vom Bremsweg und unterliegt einem gewissen Sicherheitspuffer.
\item[SenseIvtAndProtectiveField]  Diese Modul realisiert die AND--Operation zwischen dem jeweiligen
 Virtual protective Field und der Hinderniskarte. Außerdem wird der entsprechende Controlleroutput gesetzt
 sollten nach der Operation noch mehr als 3 weiße Punkte im resultierenden Bitmap vorhanden sein.
\end{description}

Die Module \lstinline{ProtectiveFieldSideward},
 \lstinline{ProtectiveFieldForward} und
 \lstinline{ProtectiveFieldCom-}\lstinline{bined} dienen wie das Modul
 \lstinline{ObstacleMap} aus der \lstinline{ObstacleExtractor}--Gruppe zur Visualisierung, wobei hier die Virtual
 protective Fields dargestellt werden können. Die Module 
 \lstinline{ProtectiveField-}\lstinline{SidewardAnd},
 \lstinline{ProtectiveFieldForwardAnd} und \lstinline{ProtectiveFieldCombinedAnd} können dazu genutzt werden die
 Bitmaps nach der AND--Operation anzuzeigen.


\chapter{Tracking und Gestenerkennung}

\todo[inline]{Schreiben}



\chapter{Integration}
\label{integration_cha}
\authorsection{\editorjulian}
\todo[inline]{Julian: Schreiben}

\section{Aufgabenstellung}
\label{aufgabenstellung_integration_sec}

Die Aufgaben im Bereich der Integration lassen sich unter folgende Punkte einordnen:

\begin{description}
\item[Code Zusammenführen]
Um das Auseinanderdriften der einzelnen Teilprojekte zu vermeiden ist ein regelmäßiges Zusammenführen 
mittels Git notwendig. Zusätzlich ist die Integration von Code der außerhalb des SegwayOmni-Projektes 
entsteht, wie etwa der Kinect-Teil, notwendig. Weiterhin muss die Kompatibilität der Schnittstellen 
zwischen den Teilprojekten gepflegt werden.

\item[Architektur]
Das bestehende hierarchische Grundkonzept muss erweitert werden um eine sinnvolle Kontrolle der 
einzelnen Funktionen des Roboters zu ermöglichen.

\item[Benutzerschnittstelle]
Zur Kontrolle des Roboters durch Benutzer muss eine Benutzerschnittstelle implementiert werden die es 
erlaubt die einzelnen Funktionen der Software zu aktivieren.

\end{description}

\section{Umsetzung}
\label{umsetzung_integration_sec}

\subsection{Architektur}
\label{integration_architektur_sec}

Die Gruppen \lstinline{SegwayOmniBehaviours} und \lstinline{SegwayOmniKinect} verfügen über 
verschiedene Funktionen und bieten wichtige Informationen über den Zustand des Roboters und 
seiner Umwelt. Um diese Informationen sinnvoll verarbeiten zu können und ein Abrufen der einzelnen 
Funktionen zu ermöglichen, ist es notwendig eine Architektur bereitzustellen die dies ermöglicht.

\subsubsection{Hierarchische Struktur}
\label{integration_umsetzung_hierarchie_sec}

In den existierenden Projekten \lstinline{Odete} und \lstinline{SegwayOmni} gab es bereits eine Hierarchie 
in den verwendeten Gruppen, bei der Schritt für Schritt die Abstraktion von der verwendeten 
Plattform bzw. Hardware erfolgte. Diese Hierachie wurde beibehalten und kann in drei Ebenen unterteilt 
werden (siehe \ref{bahnplanung_umsetzung_sec}). Auf der obersten Ebene, der Verhaltensebene, befinden sich die 
Gruppen \lstinline{SegwayOmniBehaviours} und \lstinline{SegwayOmniKinect}. Zur Kontrolle des Roboters kommt 
eine zusätzliche Ebene hinzu, die es erlaubt die einzelnen Verhalten zu aktivieren. Diese Ebene arbeitet auf 
Aufgabenbasis und aktiviert das notwendige Verhalten um eine Aufgabe auszuführen. Eine Aufgabe ist 
beispielsweise: Zeige den gespeicherten Pfad einem Menschen. Auf der Verhaltensebene muss nun die entsprechende 
Funktion aktiviert werden. In diesem Beispiel würde die
\lstinline{SegwayOmni-}\lstinline{Kinect}-Gruppe das Verhalten zum Beobachten
eines Menschen ausführen und die \lstinline{SegwayOmni-Behaviours}-Gruppe das
Verhalten zum Abfahren eines gespeicherten Pfades.

Die Funktionen des Roboters können, auf oberster Ebene, in einzelne, unabhängige Aufgaben getrennt werden. Zur Erfüllung dieser Aufgaben müssen wiederum bestimmte Verhalten ausgeführt werden.
\begin{table}[h]
	\label{tab:integration_aufgaben}
	\centering
	\begin{tabular}{|p{3.7cm}|p{3cm}|p{3.2cm}|}
	\hline
	\textbf{Aufgabenebene}	&	\multicolumn{2}{c|}{\textbf{Verhaltensebene}}\\
		&	\textbf{Bahnpl./Steur.}	&	\textbf{Kinect}	\\
	\hline
	Initialisierung	&	Initialisierung	&	Initialisierung \\
	\hline
	Warten auf Person	&	Idle	&	Warten auf Person + Warten auf Geste \\
	\hline
	Warten auf Geste	&	Idle	&	Warten auf Geste \\
	\hline
	Verfolge Person	&	Verfolge Person + Zeichne Pfad auf	&	Tracke Person + Warte auf Geste \\
	\hline
	Kehre zum Start zurück	&	Kehre zum Start zurück	&	Warte auf Geste \\
	\hline
	Zeige Pfad	& Zeige Pfad	&	Tracke Person + Warte auf Geste \\
	\hline
	\end{tabular}
	\caption{Darstellung der Aufgaben und der entsprechenden Verhalten.}
\end{table}
In einem ersten Schritt ist es sinnvoll diese Funktionen zu strukturieren, in dem jeder Aufgabe ein 
interner Zustand zugewiesen wird. Da bei diesem Projekt nicht der Dialog mit dem Roboter sondern lediglich 
die Aktivierung verschiedener Funktionen durch Gesten gewünscht ist, erfolgt der Wechsel der internen 
Zustände durch eine deterministische Zustandsmaschine. Bei komplexeren 
Aufgaben, oder in Szenarien bei der die beobachtete Umwelt nicht deterministisch aufgefasst wird, stößt 
eine einfache Zustandsmaschine an ihre Grenzen und es muss eine Planer verwendet werden.
Neben der globalen Zustandsmaschine im Modul
\lstinline{SegwayOmniStateMode-Manager} erhält jede Gruppe in der darunterliegenden Hierarchieebene (die Verhaltensebene) ebenfalls eine Zustandsmaschine; 
realisiert sind diese in den Modulen \lstinline{SegwayOmniBehavioursManager} bzw. \lstinline{SegwayOmniKinectManager}.

Mit Hilfe dieser hierarchischen Architektur und durch das Arbeiten mit Zuständen kann effizient und für 
den Benutzer stets transparent, die Robotersoftware kontrolliert werden.

\subsubsection{Zusammenspiel der Zustandsmaschinen}
\label{integration_umsetzung_zusammenspiel_sec}

Ein gängiges Konzept zur Steuerung von Robotern bildet 
eine \textit{Sense-Plan-Act-Schleife}\footnote{Dt.: Beobachten-Planen-Handeln-Schleife}, bei der 
ständig die Umwelt sensorisch erfasst und entsprechend der Beobachtung und des aktuellen Zustandes 
eine Handlung ausgeführt wird.
\begin{figure}[h]
	\label{fig:integration_sense_plan_act}
	\centering
	\includegraphics[scale=0.5]{graphics/SCHEMA-SensePlanAct.pdf}
	\caption{Schematische Darstellung des Konzeptes der Sense-Plan-Act-Schleife.}
\end{figure}
Die globale Zustandsmaschine bildet in diesem Konzept das Plan-Modul. Sie erhält von den Gruppen auf 
der Verhaltensebene Informationen, auf Basis derer die Entscheidung für einen Zustandswechsel 
getroffen wird. Ändert sich der Zustand der globalen Zustandsmaschine kann, muss aber nicht, eine 
Handlung folgen. Eine Handlung wird ausgeführt in dem die globale Zustandsmaschine, durch das 
Setzen der entsprechenden Zustände auf Verhaltensebene, ein bestimmtes Verhalten aktiviert oder 
auch ein momentan aktives Verhalten beendet.

Die \lstinline{Mca2}-Architektur bietet, durch die strikte Trennung des Informationsflusses in 
einen \lstinline{Con-trol}- und \lstinline{Sense}-Teil, eine passende Grundlage
für das zuvor beschriebene Konzept. Jede Zustandsmaschine ist in zwei Teile gespalten. Im \lstinline{Sense}-Teil 
erfolgt die Verarbeitung neuer Informationen. Dies geschieht indem, entsprechend dem vorherrschenden 
inneren Zustand und der beobachteten, externen Variable, ein Übergang zu einem neuen Zustand 
ausgeführt wird. Beispielsweise kommt die Ausführung der Aufgabe \textit{Verfolge einen Menschen} 
durch die Geste \textit{Stop} zu ihrem Ende. Im \lstinline{Control}-Teil werden hingegen 
Handlungen ausgeführt. Hier erfolgt, entsprechend dem momentanen inneren Zustand, zum Beispiel 
das Setzen eines neuen Zustandes für in der Hierarchie tiefer liegende Zustandsmaschinen. 
Ein passendes Beispiel ist der Zustand \lstinline{stopFollowing-Person} der
globalen Zustandsmaschine, welcher wiederum den Zustand \lstinline{stopFollowingPerson} im
\lstinline{SegwayOmniBehavioursManager} setzt und so das Beenden des momentanen Verhaltens initialisiert. Erst wenn die aufgezeichnete 
Trajektorie gespeichert und der Roboter gestoppt ist, erfolgt die Rückmeldung, indem die 
Zustandsmaschine in den Zustand \lstinline{idle} wechselt. Dieser neue Zustand wird 
im \lstinline{Sense}-Teil der globalen Zustandsmaschine verarbeitet und falls alle weiteten 
Bedingungen erfüllt sind wird die Aufgabe durch einen letzten Zustandswechsel, ebenfalls in 
den Zustand \lstinline{idle}, beendet (siehe Abb. \ref{fig:integration_statemachine}).


\subsection{Benutzerschnittstelle}
\label{benutzerschnittstelle_integration_cha}

Gemäß der Konzeption des Roboters erfolgt die Steuerung vollständig über Gesten. Die Erkennung der Gesten 
geschieht in der Gruppe \lstinline{SegwayOmniKinect}; die Verarbeitung der Gesten im Modul \lstinline{SegwayOmniStateModeManager}. Entsprechend der in Abschnitt \ref{integration_umsetzung_zusammenspiel_sec}
 erläuterten Systematik, entscheidet die globale Zustandsmaschine in jedem Zustand individuell, welche Handlung 
 nach der Beobachtung einer bestimmten Geste folgt. Beim ersten Test der Gestensteuerung hat sich gezeigt, 
 dass eine zuverlässige Steuerung nur durch das Erkennen einer Abfolge mehrerer Gesten möglich. 

Es trat das Problem auf, dass während des Einnehmens einer bestimmten Haltung, die entsprechend der Definition 
der Kinect-Gruppe einer Geste entsprach, regelmäßig eine nicht gewünschte Geste erkannt wurde. Um dieses Problem 
zu lösen kommt eine Steuerung zum Einsatz, bei der nur durch eine Folge von Körperhaltungen eine Geste erkannt 
wird und so eine Funktion auslöst. Bereits bei der Initialisierung einer Person ist es notwendig eine korrekte 
Folge von Haltungen einzunehmen die signalisieren dass man auch wirklich den Roboter steuern möchte. Auch beim Ausführen einer Geste um eine bestimmte Funktion des Roboters, etwa das Abfahren eines Pfades, zu aktivieren, ist es notwendig zuerst eine Anrückhaltung einzunehmen und von dieser zu einer Haltung überzugehen die letztlich die 
Funktion aktiviert.

\begin{landscape}
\begin{figure}
	\includegraphics[scale=0.43,angle=90]{graphics/SCHEMA-GlobalStateMachine-hk.pdf}
	\caption{Zustandsdiagramm der globalen Zustandsmaschine}
	\label{fig:integration_statemachine}
\end{figure}
\end{landscape}



\chapter{Test auf realem Roboter}

\todo[inline]{Steuerungs-Gruppe: Schreiben}

\section{Testumgebung}

\begin{itemize}
	\item Aufgebaute Umgebung
	\begin{itemize}
		\item FZI, oberes Stockwerk
		\item Abgegrenzter Bereich
		\item Karte erstellen zur Lokalisierung
	\end{itemize}
	\item Geplantes Bewegungs-Szenario
\end{itemize}


\section{Aufgetretene Schwierigkeiten}

\subsection{Laserscanner}

\begin{itemize}
	\item \glqq Unsichtbare Wand\grqq\ in ca. 4m Entfernung
	\begin{itemize}
		\item Scanner nicht exakt waagerecht angebracht
		\item nicht relevant, da Entfernung von 4m ausreichend
	\end{itemize}
	\item \glqq Tote Pixel\grqq
	\begin{itemize}
		\item Fehlmessungen werden mit Wert 0 und nicht MAX belegt
		\item Roboter steht immer im Hindernis -> keine Fortbewegung
		\item Optionale Zuschaltung der Module per XML
	\end{itemize}
\end{itemize}


\subsection{Stromversorgung}

\begin{itemize}
	\item Robotereigene Stromversorgung reicht nicht aus für Roboter + Kinect
	\item Netzgerät nötig
	\item Einschränkung der Reichweite
\end{itemize}

\subsection{Kinect}

\begin{itemize}
	\item verliert getrackten Menschen bei zu schneller Bewegung
	\item Verringerung der Geschwindigkeit (Rotation sowie Translation)
\end{itemize}


\subsection{Weitere Probleme}

\begin{itemize}
	\item Odometrie auf Teppichboden
	\item SW-Bugs
	\item \ldots
\end{itemize}


\chapter{Ausblick}

\section{Erreichte Ziele}
\todo[inline]{Steuerungs-Gruppe: Schreiben}

\section{Ausblick}
\todo[inline]{Kinect-Gruppe: Schreiben}

\section{Reflexion}
\todo[inline]{Anne+Andreas: Schreiben}


\pagestyle{scrplain}
\appendix

% \chapter{Analysenmethoden zur Bestimmung der Reaktionsvariablen }

\section{Gaschromatographie}

Die Bestimmung des Wasserstoffanteils mit dem GC HP-5880 A Series erfolgt durch eine Temperaturrampe. Zuerst wird die Temperatur fünf Minuten lang bei $50 ^{\circ}C$ gehalten. Danach wird bis $130 ^{\circ}C$ ($15 ^{\circ}C/min$) aufgeheizt. Der GC besaß eine gepackte Trennsäule (Porapak Q). Als Detektor stand ein Wärmeleitfähigkeitsdetektor (WLD) zu Verfügung. 

Der GC HP 6890 ist mit zwei Säulen ausgestattet, eine für die Bestimmung von Wasserstoff (80/100 Hayesep Q, 2 m lang)  mit Stickstoff als Trägergas, und eine andere Säule (60/80 Molekularsieb, 4 m lang) für CO2, Stickstoff, Sauerstoff und Kohlenwasserstoffe ($< C4$), in der Helium als Trägergas benutzt wird. Für die Konzentrationsmessung wurde zwei Detektoren verwendet: ein Wärmeleitfähigkeitsdetektor (WLD) und ein Flammenionisationsdetektor (FID), die in Serie geschaltet waren. 

In dem WLD wird die Wärmeleitfähigkeit des Trägergases (Stickstoff) mit der der gaschromatographisch getrennten Substanzen verglichen. Der Vergleich erfolgt an einer Wheatstoneschen Brücken. Solange nur Trägergas durch die Wheatstoneschen Brücken fließt, wird die Wärme vollständig abgeführt. Wenn andere Substanzen mit geringe Leitfähigkeit in die Messzelle gelangen, entsteht ein Wärmestau am Hitzdraht (Platin), der den elektrischen Wiederstand erhöht. 

Für die anderen Komponenten in der Gasphase wird ein FID benutzt, der für Verbindungen mit C-H oder C-C Bindungen eine deutlich höhere Empfindlichkeit besitzt. Die Verbrennung der Komponenten läuft über die Bildung von Radikalen, durch die die Ionisation erhöht wird. 

Die Messung erfolgt während 31 Minuten. Die Temperatur wird von $60 ^{\circ}C$ auf $220 ^{\circ}C$ erhöht. Zuerst wird nach der Probeeinspritzung die Temperatur bei $60 ^{\circ}C$ zwei Minuten lang gehalten. Dann erfolgt eine Aufheizungsrate von $5 ^{\circ}C/min$  bis eine Temperatur von $160 ^{\circ}C$ erreicht wird. Anschließend steigt die Temperatur bis $220 ^{\circ}C$ mit einer Aufheizungsrate von $15 ^{\circ}C/min$. Bei $220 ^{\circ}C$ wird die Temperatur während  5 min konstant gehalten. 

Die Gasproben werden mit Gasspritzen aus zwei Probenehmern  gezogen.  Ein erster Probenehmer befand sich direkt an der Stelle der Phasentrennung, der zweite war hinter dem Abwasserbehälter angebracht. 
Zur Kalibrierung der Gas-Chromatographen wurden täglich Gasproben mit bekannter Zusammensetzung eingespritzt. Da die Gas-Chromatographen von mehreren Arbeitsgruppen bedient wurden, war eine Ausheizung des GC während mehreren Stunden nötig, die jedes zweites Wochenende durchgeführt wurde.  Die Versuchsplanung hängte davon ab, dass die Gasproben gemessen werden konnten.

\section{Analyse zur Bestimmung des Kohlenstoffgehalts im Abwasser}

TOC-Messgerät (DC-190 Rosemount-Dohrman) 

Das Messprinzip besteht aus einer hochtemperatur-katalysierten Verbrennung der Abwasserprobe. Die Verbrennung findet in einem Quarz-Rohr statt, das  gepackt mit dem Katalysator Pt/Al2O3 ist. Die Ofentemperatur  beträgt $800 ^{\circ}C$. Durch die katalysierte unterstützte Oxidierung mit reinem Sauerstoff ($250 ml/min$), werden die Abwasserproben vollständig zu CO2 und Wasser umgewandelt. Das entstehende CO2 wird von dem Wasser in einem Kondensator-Gas-Flüssig-Trennungssystem (Kupfer-Zinn-Adsorber) gereinigt. Die Halogene werden von dem Gasprodukt getrennt. Die reine CO2-Menge wird anhand eines Infrarotdetektors ermittelt. Anorganischer Kohlenstoff (IC) wird mit einem IC  Reaktor bestimmt, der Phosphorsäure 20 \% enthält. In diesem säuerlichen Medium werden die  anorganischen Carbonate der Abwasserprobe zu CO2 umgewandelt. Das entstehende Gas wird zum Trennungssystem geleitet und mit dem Infrarotdetektor gemessen. Die Genauigkeit dieses Messprinzips lag um etwa 2 \%. 

Küvetten-Test LCK 381 (Dr. Lange GmbH, Düsseldorf, D, Messbereich 60-735 mg/l)

Die Proben müssen verdünnt werden, wenn der Messbereich überschritten werden könnte.  Zur Bestimmung des Gesamtkohlenstoff-Gehaltes (TC) wurden $200 \mu l$ der verdünnten Abwasserprobe in die mit dem Aufschlussreagenz (Natriumperoxidsulfat) gefüllte TC-Küvette (KD 381 B) pipettiert. Nach dem Verschließen mit Originaldeckel und mehrmaligem Umschwenken wurde die TC-Küvette mit einer vorbereiteten Indikatorküvette (LCK 380/381) durch einen Membran-Doppel-deckel verbunden. Gleichzeitig wurde die TiC-Küvette (KE 381 B) mit $1 ml$ der Probe befüllt. Analog zur Aufarbeitung der TC-Bestimmung wurde die TiC-Küvette mit dem Originaldeckel verschlossen und mehrfach geschüttelt. Im Anschluß daran wurde die TiC-Küvette mit der entsprechenden Indikatorküvette verbunden. Beide Küvettenkombinationen wurden gleichzeitig in einem vorgeheizten Thermostaten LT 100 (Dr. Lange GmbH) 2 h bei $100^{\circ}C$ erwärmt. Während der Erwärmung reagieren die Abwasserproben mit dem Reagenz und es entsteht CO2, das den Indikator verfärbt. Nach dem Abkühlen auf Raumtemperatur wurden die Indikatorküvetten mit dem Spektralphotometer CADAS-200 (Dr. Lange GmbH) gemessen. Die Indikatorküvetten werden außen gut gesäubert und in den Photometer eingesetzt. Das Barcode-Ettiket wird von der CADAS-200 gelesen, dadurch werden die entsprechenden Wellenlängen für die Durchführung der Messung eingestellt.   



\listoffigures
\listoftables
\bibliography{content/bibliography}

\printindex


\end{document}
